\documentclass[10pt]{article}

\usepackage[utf8]{inputenc}

\usepackage[T1]{fontenc}

\usepackage{amsmath}

\usepackage{amsfonts}

\usepackage{amssymb}

\usepackage[version=4]{mhchem}

\usepackage{stmaryrd}

\usepackage{graphicx}

\usepackage[export]{adjustbox}

\graphicspath{ {./images/} }

\usepackage{bbold}



\begin{document}

следует, что $f \in \tilde{C}^{r}$, причем аналогично (2) можно оценить $\omega\left(f^{(r)} .1 / n\right)$ через $E\left(f, T_{n-1}\right)(n=1,2, \ldots)$ (см. [4], [8], [12]). Из этих оценок, в частности, следует, что если



\[

E\left(f, T_{n-1}\right)=O(n-r-\alpha), \quad r=0,1, \ldots, 0<\alpha \leqslant 1,

\]



то $f \in \tilde{C^{r}}$ и $f^{(r)}(t)$ удовлетворяет при $0<\alpha<1$ у с л ов и ю $\Gamma$ ё ль д е p a



\[

\left|f^{(r)}\left(t^{\prime}\right)-f^{(r)}\left(t^{\prime \prime}\right)\right| \leqslant K\left|t^{\prime}-t^{\prime \prime}\right|^{\alpha},

\]



а при $\alpha=1-$ услови ю 3 и гмунда



\[

\left|f^{(r)}\left(t^{\prime}\right)-2 f^{(r)}\left(\frac{t^{\prime}+t^{\prime \prime}}{2}\right)+f^{(r)}\left(t^{\prime \prime}\right)\right| \leqslant K\left|t^{\prime}-t^{\prime \prime}\right| .

\]



Обозвачив этот класс функций через $K \tilde{H}^{r+\alpha}$, получают өго ковструктивную характеристику: $f \in K \tilde{H}^{r+\alpha}$ тогда и только тогда, когда



\[

E\left(f, T_{n-1}\right)=O(n-r-\alpha) .

\]



\includegraphics[max width=\textwidth, center]{2023_07_08_6c70d0d6a3b42e8d21a2g-1}

руема на всей оси в том и только в том случае, если



\[

\lim _{n \rightarrow \infty} n^{r} E\left(f, T_{n \rightarrow 1}\right)=0

\]



для любого $r=1,2, \ldots$



Аналогичные факты имеют место для приближения периодич функций в метрике $L_{p}[0,2 \pi](1 \leqslant p<\infty)$, а также для заданных на всей оси (не обязатөльно периодич ) функций в случае приближения их целычи функциями конечной степени (см. [7], [8]). Известны п. т. И о. т. для $f \in \bar{C}^{r}$, использующие в качестве дифференциально-разностной характеристики модуль гладкости $\omega_{k}(f, \delta)$ порядка $k=1,2, \ldots$ приближаемой Функции (или нек-рой ее производной) (см. [4], [8]).



Инаqе обстоит дело в случае приближения на конечном отрөзке. Пусть $C=C[a, b]$ - пространство непрерывных на $[a, b]$ функций с норчой



\[

\|f\|_{C}=\max _{a \leqslant t \leqslant b}|f(t)|,

\]



$C^{r}=C^{r}[a, \quad b]-$ множество $r$ раз непрерывно дифференцируемых на $[a, b]$ функций, $C^{0}=C, K H^{r+\alpha}[a, b]-$ класс функций, определяемый неравенствами (3) и (4) при $t^{\prime}, t^{\prime \prime} \in[a, b]$. Для наилучшего приблияения $E\left(f, A_{n-1} ; a, b\right)=\inf _{p \in A_{n-1}}\|f-p\|_{C}, n=1,2, \ldots$, функции $f \in C^{r}$ подпространством $A_{n-1}$ алгебраич. многочленов степени $n-1$ справедлива оценка вида (1) через модуль непрерывности функции $f^{(r)}$ на $[a, b]$, однако обращевие, аналогичное периодич. случаю (с нөравенством вида (2)), здесь возможно лишь на отрезке, лежащем внутри интервала $(a, b)$. Напр., если



\[

\begin{gathered}

E\left(f, A_{n-1} ; a, b\right) \leqslant M^{-r-\alpha}, \\

r=0,1, \ldots, \quad 0<\alpha \leqslant 1,

\end{gathered}

\]



то можно лишь утверждать, что $f$ принадлежит классу $K H^{r+\alpha}\left[a_{1}, b_{1}\right]$, определяемому неравенствами (3) (при $0<\alpha<1$ ) и (4) (при $\alpha=1$ ) лишь на отрезке $\left[a_{1}, b_{1}\right] \subset$ $\subset(a, b)$, Іричем константа $K$ зависит от $a, a_{1}, b_{1}$ и $b$ и может неограниченно увеличиваться, если $a_{1} \rightarrow a$, $b_{1} \rightarrow b$, Существуют функции, не принадлежащие классу $K H^{r+\alpha}[a, b]$, для к-рых, однако,



\[

E\left(f, A_{n-1} ; a, b\right)=O(n-r-\alpha) .

\]



Напр.,



$E\left(\sqrt{1-t^{2}}, A_{n-1} ;-1,1\right)<2 / \pi n, \quad n=1,2, \ldots$,



хотя $\sqrt{1-t^{2}} \notin K H^{\alpha}$ на $[-1,1]$ ни при каком $\alpha>1 / 2$. Оказалось, что алгебраич. многочлены могут, обеспечивая на всем отрезке $[a, b]$ ваилучший порядок при- ближения функции $f \in C$, у концов отрезка осуществлять приближение существенно лучшее (впервые этот феномен был обнаружен С. М. Никольским, см. [3]) Если, в частности, $f \in K H^{r+\alpha}[a, b]$, то при каждом $n>r$ существует многочлен $p_{n}(t) \in A_{n-1}$ такой, что



\[

\begin{aligned}

\left|f(t)-p_{n}(t)\right| \leqslant M[ & \left.\frac{1}{n} \sqrt{(t-a)(b-t)}+\frac{1}{n^{2}}\right]^{r+\alpha}, \\

& a \leqslant t \leqslant b,

\end{aligned}

\]



где константа $M$ не зависит ни от $n$, ни от $t$. Это утверждение, в отличие от (5), уже можно обратить: если для $f \in C$ существует последовательность многочленов $p_{n}(t) \in A_{n-1}$ таких, что при нек-рых $r=0,1, \ldots$ и $0<\alpha \leqslant 1$ выполнено (6), то $f \in K H^{r+\alpha}[a, b]$. Известны п. т. и о. т. для $f \in C^{r}[a, b]$, использующие модуль непрерывности и модуль гладкости (см. [4], [8]).



ПІ. т., в к-рых даются порядковые оценки погренности через дифференциально-разностные характеристики приближаемой функции, доказаны для многих конкретных методов приближения (см. [6], [8], [9]), в частности для сплайнов (наилучших и интерполяционных [10]).



Известны П. т. и о. т. для приближения в хаусдорфовой метрике (см. [13]). Здесь возникают свои особенности; в частности, характеризация классов функций через их наилучшие хаусдорфовы приближения связана не только с порядком этого приближения, но и с величиной константы в соответствующем неравенстве. О п. т. и о. т. в многомерном случае см. Приближение функций; случай многих действительных переменных.



Jum.: [1] J a c k s o n D., Über die Genauigkeit der Annäherung stetiger Funktionen dürch ganze rationale Funktionen gegebenen Grades und trigonometrische Summen gegebener Ordnung, Gött., 1911; [2] Б е p н ш т е й н С. Н., O наиллучшем прибилижении непрерывных функций посредством многочленов данной степени (1912), Собр. соч., т. 1, M., 1952, с. 11-104, [3] Н и к о Ј ь с к и й С. М., "Изв. АН СССР. Сер. матем ",

1946, т. 10 , № 4, с. $295-317 ;[4]$ Д з д д ы к В. К., Введение в теорию равномерного приближения функций полиномами, М., 1977 ; [5] К о р н е й ч у к H. П., Экстремальные задачи теории приближения, М., 1976 ; [6] Т и х о м и р о в В. М., Неноторые вопросы теории шриближений, М., 1976; [7] А х и е3 е р н., И., Лекции по теории аппроксимации, 2 изд., М., 1965;



\includegraphics[max width=\textwidth, center]{2023_07_08_6c70d0d6a3b42e8d21a2g-1(1)}

тельного переменного, М., $1960 ;[9]$ К о ро в к и н

Јинейнье операторы и теория приближений, М., 1959; [10] С т е ч к и н С. Б., С у б б о т и н Ю. Н., Сплайны в вычислительной математике, М., 1976; [11] Д а у г а в е т И. К.,

Введөние в теорию приближения функции, Л., 1977; [12] С т е чк и н С. Б., "Изв. АН СССР. Сер. матем.", 1951, т. 15, № 3 , c. 219-42; [13] С е н д о в Б л., Хаусдорфовые приближения, ШРИБЛИЯЕНИЕ ФУНКЦИИ; с п. П. Корнейчук. ПРИБЛИЯЕНИЕ ФУНКЦИИ; с л У ч а Й М н ог их дей с т ви те льны случай, когда приближаөмая функция $f$ зависит от двух и большего числа переменных:



\[

f(t)=f\left(t_{1}, t_{2}, \ldots, t_{m}\right), m \geqslant 2

\]



(см. Приближение функций). По сравнению с одномерным случаем исследование вопросов приближения функций $m(m \geqslant 2)$ пөремевных звачительно усложняется ввиду появления принципиально новых обстоятельств. связанных с многомерностью. Прежде всего это касается области, на к-рой осуществляется приближение. Просто связный компакт (в одномерном случае - отрезок) в $\mathbb{R}^{m}$ (даже на плоскости) может иметь самую разнообразную конфигурацию, и возникает необходимость классифицировать области в зависимости, напр., от гладкостных свойств их границ. Сложности появляются и при описании дифференциально-разностных свойств функций $m$ переменных. Эти свойства могут быть, вообще говоря, различными по разным направлениям, их характеризация должна учитывать как геометрию области определения, так и поведение функции при подходе к границе, так что большое значение приобретает изучение гра-





\end{document}